% BHU PYQ -- B.A./B.Sc. (Hons) Semester IV Examination 2025-2026 (NEP)
% Paper: MATMJ-42 : Vector and Tensor Analysis
% Exam: B.A./B.Sc. (Hons) Semester IV Examination, 2025-2026 (NEP)

\documentclass[11pt,a4paper]{article}
\usepackage[a4paper,top=2.5cm,bottom=2.5cm,left=2.8cm,right=2.8cm,headheight=14pt]{geometry}
\usepackage{amsmath,amssymb}
\usepackage[T1]{fontenc}\usepackage[utf8]{inputenc}
\usepackage{lmodern,microtype,fancyhdr,enumitem,hyperref,lastpage,parskip}

\hypersetup{colorlinks=true,urlcolor=black,linkcolor=black,
  pdftitle={MATMJ-42 Vector and Tensor Analysis -- BHU B.Sc. Sem IV 2025-26 (NEP)}}
\pagestyle{fancy}\fancyhf{}
\renewcommand{\headrulewidth}{0.4pt}\renewcommand{\footrulewidth}{0.4pt}
\fancyhead[L]{\small   \textit{Banaras Hindu University}}
\fancyhead[R]{\small   \textit{MATMJ-42: Vector and Tensor Analysis (2025-26 NEP)}}
\fancyfoot[C]{\small Page  \thepage\ of \pageref{LastPage}}


\newlist{parts}{enumerate}{2}
\setlist[parts,1]{label=(\alph*),leftmargin=*,itemsep=5pt,topsep=3pt}
\setlist[parts,2]{label=(\roman*),leftmargin=*,itemsep=3pt,topsep=2pt}

\newcommand{\pts}[1]{\hfill   \textit{[#1]}}


\begin{document}

\begin{center}
  {\large\bfseries BANARAS HINDU UNIVERSITY}\\[2pt]
  {\large\bfseries B.A./B.Sc. (Hons) Semester IV Examination, 2025-26 (NEP)}\\[4pt]
  {
\normalsize Subject: Mathematics \hfill Paper No. MATMJ-42 : Vector and Tensor Analysis}\\[4pt]
  \rule{0.8\linewidth}{0.5pt}\\[6pt]
  {
\normalsize Time: 2½ Hours \hfill Full Marks: 70}\\[2pt]
  {\small REF NO. Conf/Even/2025-26/65}
\end{center}
\rule{\linewidth}{0.4pt}
\smallskip



\noindent   \textit{Instructions: Answer   \textbf{five} questions including Question 1 which is compulsory. Make proper depiction of lines, curves, surfaces and volumes wherever necessary. Figures in right-hand margin indicate marks.}
\bigskip




\noindent   \textbf{Question 1.}   \textit{(Compulsory)}

\begin{parts}
  \item Make sketch of a simple curve, simple closed curve, simply connected domain and multiply connected domain.\pts{2}
  \item Find $
abla^2\left(\frac{1}{r}\right)$, where $\vec{r}$ is the position vector and $r=|\vec{r}|$.\pts{2}
  \item In which direction from the point $(2,1,-1)$ is the directional derivative of $\phi=x^2 y z^3$ maximum and find its magnitude.\pts{2}
  \item Write down the expression for the Christoffel symbols of first and second kind $[pq,r]$ and $
\begin{Bmatrix} s \\ pq \end{Bmatrix}$ respectively, where $p,q,r,s$ denotes indices.\pts{2}
  \item Prove that the contraction of the tensor $A^p_q$ is a scalar or invariant.\pts{2}
  \item Find the area of the region bounded by $y^3=x$, the $y$-axis and the lines $y=3$ and $y=6$.\pts{3}
  \item Find $
abla  \times \left(\frac{a\vec{r}}{r^3}\right)$ where $a$ is a scalar and $\vec{r}$ is the position vector and $r=|\vec{r}|$. Is this field irrotational?\pts{3}
  \item Show that $\iiint_V 
abla\phi \, dV = \iint_S \phi \hat{n} \, dS$ where $V$ is the volume enclosed by the surface $S$, $\phi$ is a scalar function and $\hat{n}$ is the unit outward normal on $S$.\pts{3}
  \item If $A_i$ and $B^i$ are the components of a covariant and contravariant vector respectively, then show that $A_i B^i$ is an invariant.\pts{3}
\end{parts}
\medskip\rule{\linewidth}{0.2pt}

\medskip


\noindent   \textbf{Question 2.}

\begin{parts}
  \item Derive an expression for $
abla \phi$ in orthogonal curvilinear coordinates $(u_1,u_2,u_3)$. Hence show that $|
abla u_p| = h_p^{-1}$, $p=1,2,3$ and find an expression for $
abla \phi$ if the transformation is $x=az, y=bx, z=cy$.\pts{6}
  \item A particle moves along the curve $\vec{r} = (t^3-4t)\hat{i} + (t^2+4t)\hat{j}$.
  
\begin{parts}
    \item Find the length of the curve it traversed after a time $t=2$.
    \item Find the component of the velocity and acceleration in the direction $\hat{i}+\hat{j}$ after a time $t=1$.\pts{6}
  \end{parts}
\end{parts}
\medskip\rule{\linewidth}{0.2pt}

\medskip


\noindent   \textbf{Question 3.}

\begin{parts}
  \item State and prove Green's theorem on a simply connected plane region $R$.\pts{6}
  \item Evaluate $\int_C \vec{F} \cdot d\vec{r}$ for the vector field $\vec{F} = -\frac{y}{x^2+y^2}\hat{i} + \frac{x}{x^2+y^2}\hat{j}$ where the region $R$ is bounded by a circle $C$ of radius $a$ centred at $(0,0)$.\pts{6}
\end{parts}
\medskip\rule{\linewidth}{0.2pt}

\medskip


\noindent   \textbf{Question 4.}

\begin{parts}
  \item Evaluate the line integral $\oint_C (yz\,dx + zx\,dy + xy\,dz)$ where $C$ is the curve $x^2+y^2=1, z=y^2$.\pts{6}
  \item Evaluate $\iiint_V (x^2+y^2+z^2)\,dx\,dy\,dz$ where $V$ is the volume of the sphere $x^2+y^2+z^2 \le 1$.\pts{6}
\end{parts}
\medskip\rule{\linewidth}{0.2pt}

\medskip


\noindent   \textbf{Question 5.}

\begin{parts}
  \item State Gauss' divergence theorem. Using this theorem, evaluate $\iint_S \vec{F} \cdot d\vec{S}$ where $S$ is the surface of the domain $x^2+y^2 = a^2$ bounded by the planes $z=0, z=b$ and $\vec{F} = x\hat{i} - y\hat{j} + z\hat{k}$.\pts{6}
  \item If $\vec{F} = 2y\hat{i} - z\hat{j} + x\hat{k}$, evaluate $\int_C \vec{F}  \times d\vec{r}$ along the curve $x=\cos t, y=\sin t, z=2\cos t$ from $t=0$ to $t=\frac{\pi}{2}$.\pts{6}
\end{parts}
\medskip\rule{\linewidth}{0.2pt}

\medskip


\noindent   \textbf{Question 6.}

\begin{parts}
  \item Consider the transformation in the curvilinear coordinate system as $x = \frac{1}{2}(u^2-v^2), y=uv, z=z$. Determine the scale factors, the unit vectors, the area elements and the volume element.\pts{7}
  \item Suppose that $T^i$ is a contravariant vector in $\mathbb{R}^2$ and that $(T^i) = (x^2, x^1)$ in the $(x^i)$ system. Calculate $(ar{T}^i)$ in the $(ar{x}^i)$ system under the change of coordinate $ar{x}^1 = (x^2)^2 
e 0, ar{x}^2 = x^1 x^2$.\pts{5}
\end{parts}
\medskip\rule{\linewidth}{0.2pt}

\medskip


\noindent   \textbf{Question 7.}

\begin{parts}
  \item Define covariant and contravariant tensors of order two. Show that any covariant tensor of order two can be expressed uniquely as the sum of a symmetric and a skew-symmetric tensor of order two.\pts{6}
  \item Consider the transformation $x=r\cos \theta, y=r\sin \theta, z=z$. Express $ds^2$ in the form $ds^2 = \sum_{i,j=1}^3 g_{ij}\,dx^i\,dx^j$. Determine the metric $g = \det(g_{ij}), i,j=1,2,3$. Hence find the volume element in the generalized coordinate system.\pts{6}
\end{parts}

\vfill\rule{\linewidth}{0.4pt}

\begin{center}\small
\href{https://bhu-exam.vercel.app/}{Click here to visit our website}\\[4pt]
\href{https://me-aryan.vercel.app/}{Click here to visit our contributor}\\[4pt]
\href{https://quashan-me.vercel.app/}{Click here to visit our contributor}
\end{center}
\end{document}
